\begin{table}[htbp]
\centering
\caption{Kebutuhan nonfungsional sistem.}
\label{tab:kebutuhan-nonfungsional}
\small
\begin{tabular}{@{}p{1.3cm}p{3cm}p{9cm}@{}}
\toprule
\textbf{Kode} & \textbf{Atribut} & \textbf{Kebutuhan Nonfungsional} \\
\midrule
KNF-01 & \emph{Deployability} & Sistem berjalan pada laptop standar tanpa GPU. \\[2pt]
KNF-02 & \emph{Interpretability} & Prediksi dapat dijelaskan (\emph{feature importance}) dan rekomendasi tertelusur ke dokumen sumber. \\[2pt]
KNF-03 & Keamanan konten & Rekomendasi \emph{grounded} pada dokumen; parameter generatif konservatif untuk menekan halusinasi angka dosis/ambang. \\[2pt]
KNF-04 & \emph{Robustness} & Tersedia \emph{graceful fallback} pada tiap lapisan (data, \emph{retrieval}, generasi). \\[2pt]
KNF-05 & Bahasa & Keluaran menggunakan Bahasa Indonesia. \\[2pt]
KNF-06 & Keamanan klinis & \emph{Human-in-the-loop}; sistem berperan sebagai pendukung keputusan, bukan pengganti penilaian klinis. \\[2pt]
KNF-07 & \emph{Maintainability} & Pemisahan \emph{domain logic} dari antarmuka. \\[2pt]
KNF-08 & Biaya/portabilitas & Memanfaatkan komponen gratis/\emph{free-tier} dan \emph{open-source}, tanpa \emph{server} khusus untuk \emph{embedding}. \\
\bottomrule
\end{tabular}
\end{table}
